\section{The terminal radial operator and its positive density}
\label{rp-section}

We prove the finite-dimensional part of the radial positivity
argument, including the identity for the actual terminal
differential operator.  All ordered products below have the
same order and shifts as the operators in the source
\cite{alpo-uniform-source}.  The proof uses exterior powers
to establish the required identities for Capelli minors.
It then replaces a further interpolation over the spectrum
of a matrix by an orthogonal decomposition using real
interpolation polynomials.  The passage from this pointwise
density to the integrated radial forms still requires the
analytic assertions specified at the end of the section.

\subsection{Ordered minors and exterior powers}

Let $\mathcal A$ be a unital complex algebra and let
$E_{ab}\in\mathcal A$, $0\leq a,b<N$, satisfy
\begin{equation}\label{rp-gl}
 [E_{ab},E_{cd}]=\delta_{bc}E_{ad}-\delta_{da}E_{cb}.
\end{equation}
For increasing lists $I=(i_0,\ldots,i_{n-1})$ and
$J=(j_0,\ldots,j_{n-1})$, with $0\leq n\leq N$, put
\begin{equation}\label{rp-minor}
 Q_{I,J}(u)=
 \sum_{\sigma\in\mathfrak S_n}\operatorname{sgn}(\sigma)
 \prod_{a=0}^{n-1}
 \bigl((u-a)\delta_{i_{\sigma(a)},j_a}
                  +E_{i_{\sigma(a)},j_a}\bigr).
\end{equation}
The factors are multiplied from left to right in increasing
$a$.  An empty product is $1$.  Write $Q_N(u)$ for the
minor with both lists equal to $(0,\ldots,N-1)$.
The identities in this subsection are algebraic, so they
also apply to differential operators acting on a common
space of smooth functions.

\begin{lemma}\label{rp-capelli}
For $P\in\operatorname{GL}_N(\mathbb C)$ and $R=P^{-1}$,
define
\[
 E'_{ab}=\sum_{c,d}P_{ca}R_{bd}E_{cd}.
\]
Then
\begin{equation}\label{rp-minor-covariance}
 Q_{I,J}(E';u)=
 \sum_{\substack{|A|=n\\|B|=n}}
       \det P_{A,I}\,\det R_{J,B}\,Q_{A,B}(E;u).
\end{equation}
Here subsets are always written in increasing order.
If $\mathcal A=\operatorname{End}_{\mathbb C}(W)$ for a
finite-dimensional Hilbert space $W$ and
$E_{ab}^*=E_{ba}$, then
\begin{equation}\label{rp-minor-adjoint}
                      Q_{I,J}(u)^*=Q_{J,I}(\overline u).
\end{equation}
In particular $Q_N(u)$ is self-adjoint for real $u$.
\end{lemma}

\begin{proof}
Let $V=\mathbb C^N$ and let $L_a(v)$ act on the $a$th
factor of $V^{\otimes n}$, with matrix entries
$v\delta_{ij}+E_{ij}$.  Operators in $\mathcal A$ act in
their displayed order.  Let $P_{ab}$ interchange auxiliary
factors $a,b$, with indices now $0,\ldots,n-1$.
Expanding one matrix entry using \eqref{rp-gl} gives
\begin{equation}\label{rp-rtt}
 ((u-v)1-P_{ab})L_a(u)L_b(v)
 =L_b(v)L_a(u)((u-v)1-P_{ab}).
\end{equation}
Indeed the coefficient identity, after subtracting the
commuting scalar terms, is
\[
 (u-v)[E_{ij},E_{kl}]
 =(u-v)(\delta_{kj}E_{il}-\delta_{il}E_{kj}).
\]
Thus this formula requires precisely \eqref{rp-gl}, with
the indicated sign.

Let $A_n$ be the antisymmetrizer and set
\[
 A_n=\frac1{n!}\sum_{\sigma}\operatorname{sgn}(\sigma)P_\sigma,
 \qquad L=L_0(u)L_1(u-1)\cdots L_{n-1}(u-n+1).
\]
We claim that
\begin{equation}\label{rp-fusion}
                           A_nL=A_nLA_n.
\end{equation}
For adjacent positions $a,a+1$ their spectral parameters
differ by one.  Formula \eqref{rp-rtt} therefore moves
$(1-P_{a,a+1})/2$ through that pair, reversing their
order.  This projector commutes with every other factor
$L_b$, and $A_n(1-P_{a,a+1})/2=A_n$.  It follows that
$A_nL(1+P_{a,a+1})=0$.  Hence $A_nLP_\sigma
=\operatorname{sgn}(\sigma)A_nL$ for every permutation,
which proves \eqref{rp-fusion}.

Use the orthonormal exterior basis
\[
 \omega_I=\frac1{\sqrt{n!}}
      \sum_{\sigma}\operatorname{sgn}(\sigma)
      e_{i_{\sigma(0)}}\otimes\cdots\otimes e_{i_{\sigma(n-1)}}.
\]
The coefficient of $\omega_I$ in $A_nL(e_{j_0}\otimes
\cdots\otimes e_{j_{n-1}})$ is $Q_{I,J}(u)/\sqrt{n!}$.
On the other hand the antisymmetrization of the input is
$\omega_J/\sqrt{n!}$.  Equation \eqref{rp-fusion} proves
that the matrix entries of $A_nLA_n$ on $\bigwedge^nV$
are exactly $Q_{I,J}(u)$, with no additional factorial.

The auxiliary matrix $E'$ is $P^tE(P^t)^{-1}$.
Conjugate every auxiliary factor by $P^t$.
The antisymmetrizer commutes with this simultaneous change
of basis.  The entries of $\bigwedge^nP^t$ and
$\bigwedge^nR^t$ are respectively $\det P_{A,I}$ and
$\det R_{J,B}$.  This proves
\eqref{rp-minor-covariance}.

For the adjoint assertion first take $u$ real.  Each
$L_a(u-a)$ is self-adjoint because $E_{ij}^*=E_{ji}$.
Inside an expression $A_n(\cdots)A_n$ one may interchange
two neighboring factors whose labels are $a<b$.
To see this, use \eqref{rp-rtt} with parameter difference
$b-a$.  Its factor $(b-a)1-P_{ab}$ commutes with all the
other $L_c$ and acts on either exterior space as the
nonzero scalar $b-a+1$.  Move this factor across the
expression and cancel that scalar.  Reversing the ordered
list $0,\ldots,n-1$ by adjacent interchanges only exchanges
pairs with $a<b$ in this direction.  Consequently
\[
 A_nL_0(u)\cdots L_{n-1}(u-n+1)A_n
 =
 A_nL_{n-1}(u-n+1)\cdots L_0(u)A_n.
\]
The right side is the adjoint of the left side.
Comparison of exterior matrix entries proves
\eqref{rp-minor-adjoint} for real $u$.  Both sides are
polynomials in $\overline u$, so equality for all real
$u$ proves the general assertion.  For $n=0$ all these
identities reduce to $1=1$.
\end{proof}

\begin{corollary}\label{rp-central-character}
The full determinant $Q_N(u)$ commutes with every $E_{ab}$.
Suppose a module for these operators contains a vector $v$
such that
\[
 E_{ab}v=0\quad(a<b),\qquad E_{aa}v=\lambda_a v.
\]
Then
\begin{equation}\label{rp-highest-eigenvalue}
                 Q_N(u)v=\prod_{a=0}^{N-1}(u-a+\lambda_a)v.
\end{equation}
The same scalar identity holds on the submodule generated
by $v$.
\end{corollary}

\begin{proof}
For $n=N$, \eqref{rp-minor-covariance} says $Q_N(E';u)
=Q_N(E;u)$, because $\det P\det P^{-1}=1$.
Take $P(t)=\exp(tX)$ and differentiate at $0$.
The derivative of $E'_{ab}$ is
$[\sum_{c,d}X_{cd}E_{cd},E_{ab}]$ by \eqref{rp-gl}.
The product rule therefore gives
$[\sum X_{cd}E_{cd},Q_N(u)]=0$, proving centrality.
This argument may equally be read as an identity of
formal power series, so no topology on $\mathcal A$ is
needed.

In a summand of \eqref{rp-minor} belonging to a nonidentity
permutation, let $a$ be its largest nonfixed column.
Then $\sigma(a)<a$.  All factors to its right are diagonal
and act on $v$ by scalars; its factor is $E_{\sigma(a),a}$
and annihilates $v$.  Only the identity permutation
remains, yielding \eqref{rp-highest-eigenvalue}.
Centrality extends the identity from $v$ to every vector
obtained by applying the $E_{ab}$, and thence to their
linear span.
\end{proof}

\begin{corollary}\label{rp-negative-convention}
For the alternative convention
\[
 B_N(u)=\sum_{\sigma}\operatorname{sgn}(\sigma)
           \prod_{a=0}^{N-1}
             \bigl((u+a)\delta_{\sigma(a),a}-E_{\sigma(a),a}\bigr)
\]
one has $B_N(u)=(-1)^NQ_N(-u)$.  It is central,
self-adjoint for real $u$ under the preceding adjoint
hypothesis, and acts on the highest-weight submodule
above by $\prod_a(u+a-\lambda_a)$.
\end{corollary}

\begin{proof}
Factor $-1$ out of every factor in the ordered product.
The other conclusions follow from the two preceding
results.
\end{proof}

\subsection{The complementary contraction}

For $N\geq2$, vectors $a,a'\in\mathbb C^N$ and covectors
$\beta,\beta'\in(\mathbb C^N)^*$, define
\begin{align}
 \mathcal C_E(u;a,a',\beta,\beta')
 ={}&\sum_{\substack{i<i'\\j<j'}}
       (-1)^{i+i'+j+j'}
       (\beta_i\beta'_{i'}-\beta_{i'}\beta'_i)
       (a_ja'_{j'}-a_{j'}a'_j)\notag\\
 &\hspace{32mm}\cdot Q_{\{i,i'\}^{c},\{j,j'\}^{c}}(u).
 \label{rp-contraction}
\end{align}
All coordinate factors in this formula are scalars placed
to the left of the differential operators, when these
are the chosen algebra $\mathcal A$.

\begin{lemma}\label{rp-adapted-contraction}
Let $P=(a,a',b_1,\ldots,b_{N-2})$ be an invertible matrix.
Suppose $\beta(b_i)=\beta'(b_i)=0$ for all $i$, and let
$E'$ be its change of generators from Lemma~\ref{rp-capelli}.
Then
\begin{equation}\label{rp-contraction-adapted}
 \mathcal C_E(u;a,a',\beta,\beta')
 =\bigl(\beta(a)\beta'(a')-\beta(a')\beta'(a)\bigr)
     Q_{\{2,\ldots,N-1\},\{2,\ldots,N-1\}}(E';u).
\end{equation}
\end{lemma}

\begin{proof}
We first verify covariance of \eqref{rp-contraction}.
For two-element subsets $I,J$ put
\[
       C_{I,J}(E)=(-1)^{\sum I+\sum J}Q_{I^c,J^c}(E;u).
\]
Complementary minors of an invertible scalar matrix obey
\[
 \det P_{A^c,I^c}
   =(-1)^{\sum A+\sum I}\det P\,\det(P^{-1})_{I,A}.
\]
For completeness this is the coefficient identity obtained
by applying $\bigwedge^NP$ to
$e_A\wedge e_{A^c}$ and comparing the wedge pairing
\[
 \bigwedge\nolimits^2\mathbb C^N\times
 \bigwedge\nolimits^{N-2}\mathbb C^N
                   \longrightarrow\bigwedge\nolimits^N\mathbb C^N.
\]
The sign of $e_A\wedge e_{A^c}$ is
$(-1)^{\sum A-1}$; the two constant signs cancel.
Equivalently, expanding that wedge pairing is the usual
Laplace expansion, which proves the displayed identity
without any assumption on the complementary minor.

Apply \eqref{rp-minor-covariance} and this identity twice.
Since $\det P\det P^{-1}=1$, one obtains
\[
 C(E')=(\bigwedge\nolimits^2P^{-1})\,C(E)\,
                                     (\bigwedge\nolimits^2P).
\]
The vector $a\wedge a'$ changes by $\bigwedge^2P^{-1}$,
and the covector $\beta\wedge\beta'$ by $\bigwedge^2P$.
Their contraction with $C$ is consequently unchanged.
In the basis formed by the columns of $P$ the two vectors
are $e_0,e_1$, and both covectors vanish in positions
$2,\ldots,N-1$.  Only $I=J=\{0,1\}$ survives in
\eqref{rp-contraction}; its sign is $+1$.  This gives
\eqref{rp-contraction-adapted}, including the case $N=2$,
whose remaining determinant is the empty determinant $1$.
\end{proof}

\subsection{The actual terminal differential operator}

Let $V$ be a complex Hermitian space of signature $(k,1)$,
$k\geq1$, with its Hermitian form $h$ linear in the first
argument.  Let $z,z'\in V$ satisfy
$h(z,z)=-T$, $h(z',z')=-T'$ with $T,T'>0$.
For $g\in U(V,h)$ put
\[
 y=g^{-1}z,\qquad x=z',\qquad
 \zeta=h(gz',z)=h(x,y),\qquad
 G(g)=|\zeta|^2-TT'.
\]
We work at points where $G(g)>0$.  On the terminal group
of the source, the previously fixed isotropic vectors
have been removed; the remaining space is exactly this
$(k+1)$-dimensional space $V$.

For $A\in\operatorname{End}_{\mathbb C}(V)$ write
$A^\dagger$ for its $h$-adjoint and define
\[
 A_1=\tfrac12(A-A^\dagger),\qquad
 A_2=\tfrac1{2i}(A+A^\dagger).
\]
Both belong to the real Lie algebra $\mathfrak u(V,h)$,
and $A=A_1+iA_2$.  If $F$ is a smooth function on an
open part of $U(V,h)$, define
\[
 R_XF(g)=\left.\frac{d}{dt}\right|_{t=0}F(g\exp(tX)),
 \qquad E_AF=R_{A_1}F+iR_{A_2}F .
\]
In a fixed basis of $V$ set $E_{ab}=E_{e_a e_b^*}$,
where $e_b^*$ is the coordinate covector, not the
Hermitian dual.  The commutator formula for
left-invariant vector fields gives \eqref{rp-gl}.
It also follows directly by differentiating the group
commutator
$\exp(tX)\exp(sY)\exp(-tX)\exp(-sY)$ at $(0,0)$
and extending complex linearly.

At the fixed point $g$, put
$\beta(v)=h(v,y)$ and $\beta'(v)=h(v,x)$.
Define the terminal $(2,2)$ entry by
\begin{equation}\label{rp-actual-terminal}
 (\mathcal T_{22}(Y)F)(g)
 =-\bigl(\mathcal C_E(Y+(k-2)/2;y,x,\beta,\beta')F\bigr)(g).
\end{equation}
The point-dependent scalar coordinates of $y,\beta$
multiply the operators on the left; they are not
differentiated in this definition.
This is the terminal entry of the original two-by-two
operator: its full spectral shift is $Y+k/2$, and the
complementary contraction uses that shift minus one.

\begin{proposition}\label{rp-terminal-identity}
Choose an $h$-orthonormal basis $b_1,\ldots,b_{k-1}$
of $\langle y,x\rangle^\perp$, and put
$U_{ab}(v)=h(v,b_b)b_a$.  Let
$Q_{k-1}(E_U;u)$ be the ordered full determinant of
the operators $E_{U_{ab}}$.  Then
\begin{equation}\label{rp-terminal-action}
 (\mathcal T_{22}(Y)F)(g)
       =G(g)\bigl(Q_{k-1}(E_U;Y+(k-2)/2)F\bigr)(g).
\end{equation}
The basis in this formula is held fixed while the
differential operators are applied.
\end{proposition}

\begin{proof}
The Gram matrix of $y,x$ is
$\left(\begin{smallmatrix}-T&\overline\zeta\\
\zeta&-T'\end{smallmatrix}\right)$, with determinant
$TT'-|\zeta|^2=-G(g)<0$.  It is nondegenerate of
signature $(1,1)$; its orthogonal complement is positive
definite of dimension $k-1$.  Gram--Schmidt therefore
gives the stated basis, including the empty basis when
$k=1$.

Use the basis $P=(y,x,b_1,\ldots,b_{k-1})$ of $V$.
The covectors $\beta,\beta'$ vanish on its last $k-1$
columns, and
\[
 \beta(y)\beta'(x)-\beta(x)\beta'(y)
                    =TT'-|\zeta|^2=-G(g).
\]
For a last-column basis index $b$, the corresponding
row of $P^{-1}$ is $v\mapsto h(v,b_b)$.
Consequently the transformed generator $E'_{ab}$ in
the last $(k-1)\times(k-1)$ block is $E_{U_{ab}}$:
the rank-one matrix underlying it is
$P_{\cdot,a}(P^{-1})_{b,\cdot}=b_a h(\,\cdot\,,b_b)$.
Lemma~\ref{rp-adapted-contraction}, followed by the
minus sign in \eqref{rp-actual-terminal}, proves
\eqref{rp-terminal-action}.
\end{proof}

Let $K_{z'}$ be the stabilizer of $z'$ in $U(V,h)$.
It is the compact unitary group of the positive space
$(z')^\perp$.  Let $W,W'$ be finite-dimensional complex
Hilbert spaces and let $\pi':K_{z'}\to U(W)$ be a
continuous representation.  Suppose that
$F:U(V,h)\supset D\to\operatorname{Hom}_{\mathbb C}(W,W')$
is smooth and satisfies
\begin{equation}\label{rp-equivariance}
                     F(gh)=F(g)\pi'(h)
\end{equation}
whenever $h\in K_{z'}$ and the relevant points lie in $D$.
It suffices that this holds in a neighborhood of each
point where a derivative is taken.

\begin{lemma}\label{rp-unitary-derivative}
A continuous finite-dimensional unitary representation
of a compact matrix Lie group is smooth.  Its
complexified derivative is a Lie algebra representation.
For the matrices $U_{ab}$ above, its operators
$D_{ab}=d\pi'_{\mathbb C}(U_{ab})$ satisfy
\[
 [D_{ab},D_{cd}]=\delta_{bc}D_{ad}-\delta_{da}D_{cb},
                      \qquad D_{ab}^*=D_{ba}.
\]
\end{lemma}

\begin{proof}
Here is a direct smoothness argument.  Normalize Haar
measure on the compact group to have mass one.  For a
smooth nonnegative function $f$ of integral one supported
near the identity, the operator
$\pi'(f)=\int f(h)\pi'(h)\,dh$ is arbitrarily close to
the identity, by continuity, and hence invertible.
For every $v$ the function
\[
 \pi'(g)\pi'(f)v=\int f(g^{-1}h)\pi'(h)v\,dh
\]
is smooth in $g$, by differentiation under the compact
integral.  Since $\pi'(f)$ is onto, every orbit map is
smooth.  Such functions $f$ exist by a smooth coordinate
chart and a nonnegative smooth bump function.  Haar
measure here may be obtained from the invariant
Riemannian volume on the compact unitary group.

Differentiating the representation identity and its
group commutator gives a real Lie algebra homomorphism;
its complex-linear extension has the same property.
Differentiating preservation of the inner product gives
$d\pi'(X)^*=-d\pi'(X)$ for a skew-Hermitian $X$.
Since $U_{ab}^\dagger=U_{ba}$, the decomposition into
the two skew-Hermitian parts above gives
$D_{ab}^*=D_{ba}$.  Finally
$U_{ab}U_{cd}=\delta_{bc}U_{ad}$ by orthonormality,
which gives the displayed commutator.
\end{proof}

\begin{corollary}\label{rp-finite-terminal}
For a fixed $g$ and its fixed basis $(b_a)$ put
\[
                 C_g(Y)=Q_{k-1}(D;Y+(k-2)/2).
\]
For every smooth equivariant $F$ as above,
\begin{equation}\label{rp-evaluation-transfer}
             G(g)^{-1}(\mathcal T_{22}(Y)F)(g)=F(g)C_g(Y).
\end{equation}
The operator $C_g(Y)$ on $W$ is self-adjoint for real
$Y$.  No choice of a smooth family of bases in $g$ is
needed for either assertion.
\end{corollary}

\begin{proof}
Every $U_{ab}$ annihilates $z'$; its two real Lie algebra
parts lie in the Lie algebra of $K_{z'}$.
Differentiating \eqref{rp-equivariance} gives
$E_{U_{ab}}F(g)=F(g)D_{ab}$.  Repeated differentiation
preserves the order: for example
$E_{U_{ab}}E_{U_{cd}}F(g)=F(g)D_{ab}D_{cd}$.
Thus every ordered minor transfers to the corresponding
minor in $D$.  Proposition~\ref{rp-terminal-identity}
proves \eqref{rp-evaluation-transfer}, and
Lemma~\ref{rp-capelli} proves self-adjointness.
\end{proof}

\subsection{A pointwise orthogonal decomposition}

We state explicitly the algebraic part of the normal
form that the positivity argument uses.  This does not
assume that the algebraic involution on differential
operators is already an adjoint for an integrated
radial form.

\begin{proposition}\label{rp-idempotent-positive}
Let $\mathcal S$ be a nonempty finite set and let
$p_\rho(Y)\in\mathbb R[Y]$, $\rho\in\mathcal S$, be
pairwise distinct.  Let $\mathcal V$ be a complex vector
space, let
$\kappa:\mathbb C^{\mathcal S}\to
\operatorname{End}_{\mathbb C}(\mathcal V)$ be a unital
algebra homomorphism, and let
$\operatorname{ev}:\mathcal V\to\operatorname{Hom}(W,W')$
be linear.  Suppose that for each real $Y$ there is a
self-adjoint $C(Y)\in\operatorname{End}(W)$ such that
\begin{equation}\label{rp-normal-evaluation}
 \operatorname{ev}\bigl(\kappa((p_\rho(Y))_\rho)\Phi\bigr)
                         =(\operatorname{ev}\Phi)C(Y)
                  \qquad(\Phi\in\mathcal V).
\end{equation}
Write $P_\rho=\kappa(1_{\{\rho\}})$ and use the
Hilbert--Schmidt form linear in the first argument,
$\langle A,B\rangle_{\rm HS}=\operatorname{tr}(B^*A)$.
Then the maps $\operatorname{ev}(P_\rho\Phi)$ are
pairwise orthogonal and
\begin{equation}\label{rp-positive-decomposition}
 \left\langle\operatorname{ev}(\kappa(f)\Phi),
                    \operatorname{ev}\Phi\right\rangle_{\rm HS}
       =\sum_{\rho\in\mathcal S}
              f(\rho)\|\operatorname{ev}(P_\rho\Phi)\|_{\rm HS}^2
                  \qquad(f\in\mathbb C^{\mathcal S}).
\end{equation}
In particular it is nonnegative for real $f\geq0$,
and is bounded between
$\min f\,\|\operatorname{ev}\Phi\|_{\rm HS}^2$ and
$\max f\,\|\operatorname{ev}\Phi\|_{\rm HS}^2$.
\end{proposition}

\begin{proof}
Choose a real $Y_0$ at which all $p_\rho(Y_0)$ are
distinct.  Such a point exists because each nonzero
difference polynomial has finitely many roots.
Put $\xi_\rho=p_\rho(Y_0)$ and form the real
interpolation polynomials
\[
 \ell_\rho(X)=
       \prod_{\sigma\ne\rho}\frac{X-\xi_\sigma}{\xi_\rho-\xi_\sigma}.
\]
Since $\kappa$ is an algebra homomorphism,
$P_\rho=\ell_\rho(\kappa((\xi_\sigma)_\sigma))$.
Induction on the degree of a polynomial in
\eqref{rp-normal-evaluation} therefore gives
\[
       \operatorname{ev}(P_\rho\Phi)
          =(\operatorname{ev}\Phi)\ell_\rho(C(Y_0)).
\]
The matrix $\ell_\rho(C(Y_0))$ is self-adjoint.
For any $\Phi,\Psi\in\mathcal V$, cyclicity of the
finite-dimensional trace gives
\[
 \langle\operatorname{ev}(P_\rho\Phi),\operatorname{ev}\Psi\rangle_{\rm HS}
 =\langle\operatorname{ev}\Phi,\operatorname{ev}(P_\rho\Psi)\rangle_{\rm HS}.
\]
On $\mathcal V$ the $P_\rho$ are mutually annihilating
idempotents whose sum is the identity.  Apply this last
identity to $P_\rho\Phi$ and $P_\sigma\Phi$ to see that
their evaluations are orthogonal for $\rho\ne\sigma$.
Their evaluations sum to $\operatorname{ev}\Phi$.
Finally $\kappa(f)=\sum_\rho f(\rho)P_\rho$, which gives
\eqref{rp-positive-decomposition} and its bounds.
No assertion about the spectrum of $C(Y_0)$ on vectors
outside these evaluated functions is needed.
\end{proof}

\subsection{The actual interlacing parameters and signs}

Here are the precise scalar parameters used in this
application.  Let $k\geq1$, $T,T'>0$, and let
$u_1,\ldots,u_k,u'_1,\ldots,u'_k$ be real numbers with
\begin{equation}\label{rp-interlacing}
 \max(u_{i+1},u'_{i+1})+1
              \leq\min(u_i,u'_i)\qquad(1\leq i<k).
\end{equation}
The source also requires these numbers to lie in a
common coset of $\mathbb Z$; the following arguments
do not require that extra condition.
Set
\[
 a_i=\max(u_{i+1},u'_{i+1})+\tfrac12,\qquad
 b_i=\min(u_i,u'_i)-\tfrac12,
\]
and let $\mathcal S$ be the product of the finite sets
$\{a_i,a_i+1,\ldots,a_i+\lfloor b_i-a_i\rfloor\}$,
$1\leq i<k$.  Thus a member $\rho$ is a tuple
$(\rho_1,\ldots,\rho_{k-1})$ with $\rho_i\in[a_i,b_i]$.
The set is nonempty; for $k=1$ it consists of the empty
tuple.  Put
\begin{equation}\label{rp-spectral-polynomials}
 p_\rho(Y)=\prod_{i=1}^{k-1}(Y+\rho_i),\qquad
 q(Y)=-TT'\prod_{i=1}^{k}(Y+\tfrac12+u_i)
                                  (Y+\tfrac12+u'_i).
\end{equation}

\begin{lemma}\label{rp-parameter-signs}
The polynomials $p_\rho$ are pairwise distinct.
Fix $1\leq l\leq k$ and a real $Y$ with
\begin{equation}\label{rp-cut-margin}
          \min(u_l,u'_l)+2\leq-Y\leq\max(u_l,u'_l)-2.
\end{equation}
Then, for every $\rho\in\mathcal S$,
\[
 (-1)^{k-l}p_\rho(Y+\delta)>0\quad(-1\leq\delta\leq1),
               \qquad q(Y)>0,\quad q(Y-1)>0.
\]
\end{lemma}

\begin{proof}
For $i<k-1$,
$\rho_i\geq\max(u_{i+1},u'_{i+1})+1/2$ and
$\rho_{i+1}\leq\min(u_{i+1},u'_{i+1})-1/2$.
Hence $\rho_i-\rho_{i+1}\geq1$.
Equality of the monic polynomials $p_\rho,p_\sigma$
gives equality of their multisets of real roots, and
the decreasing order gives $\rho=\sigma$.

For $i<l$, \eqref{rp-interlacing} implies
$\rho_i\geq\max(u_l,u'_l)+1/2$, so
$Y+\rho_i\geq5/2$.
For $i\geq l$ it gives
$\rho_i\leq\min(u_l,u'_l)-1/2$, so
$Y+\rho_i\leq-5/2$.
These signs persist after a shift of absolute value at
most one.  There are exactly $k-l$ negative factors,
proving the first assertion, also for the empty tuple.

For $i<l$, both $u_i,u'_i$ are at least
$\max(u_l,u'_l)+1$; for $i>l$ both are at most
$\min(u_l,u'_l)-1$.  Thus the two factors belonging to
each $i\ne l$ in $q(Y)$ and $q(Y-1)$ have the same,
nonzero sign.  At $i=l$, one entry is at most $-Y-2$
and the other is at least $-Y+2$ by
\eqref{rp-cut-margin}.  The two factors have opposite
signs both after adding $1/2$ and after subtracting
$1/2$.  The product in \eqref{rp-spectral-polynomials}
is negative, and its factor $-TT'$ makes it positive.
\end{proof}

For clarity we give the exact recursion for the scalar
diagonal coefficients in the source.  Fix a finite set
of positions $S$ and a list $w$ of triples
$\ell=(\varepsilon,a,y)$, where
$\varepsilon\in\{+,-\}$, $a\in S$, and $y\in\mathbb R$.
Positions are distinct and any two parameters have
distance greater than two.  Fix
$e:S\to\{0,1\}$.  For a parameter $Y$ set
\[
 \phi_\ell(Y)=
 \begin{cases}(Y-y)/(Y-y-1),&\varepsilon=+,\\
               (y-Y)/(y-Y-1),&\varepsilon=-,
 \end{cases}
 \qquad
 \eta_\ell(Y)=
 \begin{cases}(Y-y+1)/(Y-y),&\varepsilon=+,\\
               (y-Y)/(y-Y-1),&\varepsilon=-.
 \end{cases}
\]
Define $q_\varnothing(Y)=q(Y)$,
$x_{\rho,\varnothing,e}(Y)=(-1)^{k-l}p_\rho(Y)$
and $c_{\rho,\varnothing,e}=1$.  On prepending a triple
$\ell=(\varepsilon,a,y)$, set
\begin{align*}
 q_{\ell w}(Y)&=q_w(Y)\eta_\ell(Y),\\
 x_{\rho,\ell w,e}(Y)&=
       x_{\rho,w,e}(Y)
           \begin{cases}\phi_\ell(Y),&e(a)=0,\\1,&e(a)=1,\end{cases}\\
 c_{\rho,\ell w,e}&=c_{\rho,w,e}
   \begin{cases}
      q_w(y)/x_{\rho,w,e}(y+1),&\varepsilon=+,\,e(a)=0,\\
      q_w(y-1)/x_{\rho,w,e}(y-1),&\varepsilon=-,\,e(a)=0,\\
      x_{\rho,w,e}(y),&e(a)=1 .
   \end{cases}
\end{align*}

\begin{lemma}\label{rp-recursive-positive}
If every parameter $y$ in the list satisfies
\eqref{rp-cut-margin}, then all the coefficients
$c_{\rho,w,e}$ are strictly positive, and every
denominator appearing in their definition is nonzero.
\end{lemma}

\begin{proof}
If $|s|>1$, all of $s/(s-1)$ and $(s+1)/s$ are
positive.  The same statement with $s$ replaced by
$-s$ proves that $\phi_\ell(Y)$ and $\eta_\ell(Y)$
are positive whenever $|Y-y|>1$.
When a new parameter $y$ is inserted, each previously
inserted parameter $y'$ satisfies $|y-y'|>2$.
Thus $|y+\delta-y'|>1$ for $\delta=-1,0,1$.
Lemma~\ref{rp-parameter-signs} and the product formulas
give
$x_{\rho,w,e}(y),x_{\rho,w,e}(y-1),
x_{\rho,w,e}(y+1),q_w(y),q_w(y-1)>0$.
Every case of the recursion for $c$ is consequently a
positive product or quotient.  Induction on the length
proves both assertions.
\end{proof}

\begin{corollary}\label{rp-terminal-density-positive}
Take the actual terminal evaluation
$\operatorname{ev}_{c,g}\Phi=\Phi(c,g)$, with
$G(g)>0$, on the space of smooth radial functions.
Suppose that its algebraic normal form supplies a
unital algebra homomorphism $\kappa$ with
\begin{equation}\label{rp-required-normal-form}
 \kappa((p_\rho(Y))_\rho)=G^{-1}\mathcal T_{22}(Y)
                         \quad\text{on that function space}.
\end{equation}
Then \eqref{rp-positive-decomposition} holds at every
such point, with the actual parameters
\eqref{rp-spectral-polynomials}.

In particular, let $\Psi_e$ be any family of functions
in that space, indexed by $e:S\to\{0,1\}$, and let
$h$ be the length of a list satisfying the hypotheses
of Lemma~\ref{rp-recursive-positive}.  The terminal
density
\begin{equation}\label{rp-terminal-density}
 \mathcal D(c,g)=
  \left(\frac{|\zeta(g)|^2}{G(g)}\right)^h
  \sum_e\left\langle
    \operatorname{ev}_{c,g}
       \bigl(\kappa((c_{\rho,w,e})_\rho)\Psi_e\bigr),
            \operatorname{ev}_{c,g}\Psi_e\right\rangle_{\rm HS}
\end{equation}
is nonnegative.  The choice $\Psi_e=(F_w\Phi)_e$
is exactly the terminal density in the source.
If the radial functions are extended by zero off their
prescribed domain, the same conclusion holds there.
\end{corollary}

\begin{proof}
Corollary~\ref{rp-finite-terminal} gives
\eqref{rp-normal-evaluation} from
\eqref{rp-required-normal-form}.
Lemma~\ref{rp-parameter-signs} supplies distinct
polynomials, so Proposition~\ref{rp-idempotent-positive}
applies.  Each coefficient $c_{\rho,w,e}$ is positive
by Lemma~\ref{rp-recursive-positive}.
Also $|\zeta|^2=G+TT'>0$.  Every summand in
\eqref{rp-terminal-density} is therefore nonnegative.
Off the radial domain both evaluated functions are
zero, by the specified extension.
\end{proof}

\subsection{Scope and checks on the hypotheses}

The terminal identity
\eqref{rp-terminal-action} and its compact
representation interpretation
\eqref{rp-evaluation-transfer} are proved here for
the actual operators, with their exact shifts and signs.
For the subsequent positivity step the only normal-form
property used is \eqref{rp-required-normal-form};
the relations involving the additional raising and
lowering operators are unnecessary for this step.
The source derives this property as the $(2,2)$ entry
of a larger normal form.  Section~\ref{rd-section}
constructs exactly the required property directly
from compact branching.  The orthogonal decomposition
also strengthens the source's positivity conclusion to
the exact sum of squares \eqref{rp-positive-decomposition}
and its minimum and maximum bounds.

This pointwise result does not yet identify the
integrated radial form with the integral of
\eqref{rp-terminal-density}.  That identification needs
the adjoint relations on the actual class of functions,
integrability after the successive differential
operations, and the successive one-variable integral
identities with their boundary terms.  The wave
identities and the common positive majorants used in
Section~\ref{ca-section} consequently remain subject to
those analytic steps.

Two simple checks explain essential hypotheses without
giving counterexamples to the source.  A scalar operator
is not positive merely because it is self-adjoint:
for $N=1$, $E_{00}=0$, the determinant is $Q_1(u)=u$,
which is negative for real $u<0$.
The positivity above uses the positive interlacing
weights, not positivity of an arbitrary Capelli
polynomial.  Likewise replacing the strict separation
$|y-y'|>2$ by merely distinct parameters can make the
recursion undefined: with two plus parameters $y'=0$
and $y=1$, the factor $\phi_{\ell'}(y)=y/(y-1)$
has a pole.  The interpolation sets in
Section~\ref{cg-section} satisfy the stronger separation,
and Lemma~\ref{rp-parameter-signs} supplies all the
actual sign conditions.


